\section{本章小结}
\label{sec:graph-spectral-partition-summary}

组合拉普拉斯的第二小特征值$\lambda_2$衡量连通图连得有多紧。它不低于点连通度和边连通度（对简单无权非完全图），Fiedler向量的非负一侧在连通图上保持连通。但$\lambda_2$会随边权整体放大而增大，跨图比较时需要消除尺度。

归一化拉普拉斯的第二小特征值$\mu_2$用度加权长度替代普通长度。Cheeger不等式把$\mu_2$与最小导通率$\phi_*$从两侧绑在一起：下界$\mu_2/2\leq\phi_*$说明图中不存在太稀疏的平衡割，上界$\phi_*\leq\sqrt{2\mu_2}$提供了一个构造性的搜索——按特征方向排序并扫描阈值，其中至少一个割满足导通率保证。超立方体实现了下界等号，长环展示了上界中平方根尺度的必要性。

实际分组只给顶点两种值。RatioCut用顶点数计量组规模，松弛为$\lambda_2$的变分问题。NCut用度的总和计量，松弛为$\mu_2$的变分问题。两者的流程相同：求特征向量、排序、逐阈值扫描、按原目标选出最优候选。松弛值是离散最优值的下界，扫描候选给出上界，两者合在一起为分组质量提供可核验的区间。

这些步骤反复用到拉普拉斯的二次型。下一章研究怎样用更少的边同时近似保留所有顶点方向上的能量，并将这种稀疏近似用于方程求解。
